\chapter{力量}

身体力量是生命中最重要的东西。不管我们想不想，这都是真理。虽然随着人类历史的发展，身体力量对我们的日常生活来说越来越不紧要，但对我们生命的重要性却从未减弱。相比我们拥有的其他任何东西，我们的力量都在更大程度上决定了我们生存的质量和长度。以前，我们的身体力量决定了我们可以获得多少食物，以及我们的居所有多么温暖、干燥；现在，在人类通过文明的积淀所创造的新环境中，力量仅仅决定了我们的身体在其中会运转得怎样。但我们仍然是动物——归根结底，身体的存在才是唯一重要的。一个人在身体孱弱时不会比身体强壮时活得开心。对于更看重智慧和精神的那些人来说，这个事实可能难以接受。那么，看看他们在深蹲力量提升了之后的情况吧。

当我们的文化本质改变时，我们与身体活动的关系也会随之发生变化。我们之前生活在一个简单的世界中，身体的强壮作为一种素质与自身的生存紧密相连。我们曾经很好地适应了这种生存方式，因为我们别无选择。那些身体足够强壮、能够生存下来的人继续活着并且保持着强壮。这种生存方式塑造了我们的生理基础，并一直跟随我们直到现在。而就在不久前，人类社会才出现了“劳动分工”这种组织生产的新方式，所以我们的基因没有时间去再次适应。而且大多数人已经没有必要去独立地争取个人生存的权利了，身体的活动能力也被看作是可有可无的。从直接的必要性来讲，确实是这样；但我们经历了数百万年的进化所形成的强健的身体，不会因为办公桌的发明而消失。

不管是否喜欢，我们仍然拥有可以变得强壮的肌肉、骨骼、肌腱和神经，这些来之不易的宝贝需要我们的正视。这些东西经历了那么久的时间才得以形成，不能就这样被忽视了——而我们正在忽视它们，我们会自食其果的。它们是我们的肉体存在不可或缺的组成部分；它们的质量取决于我们有意识的、有目标的努力——努力提供它们需要的刺激，让它们保持在对它们来说正常的状态中。而锻炼就是这个刺激。

除了对运动表现的考虑，锻炼还是一种能够使我们的身体回归到原初设计状态的刺激。离开高强度的身体活动，人类的身体就是不正常的。锻炼不是我们针对某一个问题而找到的解决方法——而是我们无论如何都必须做的一件事，一件如果我们不做就会出现很多问题的事情。我们必须通过锻炼来复制出我们的生理功能曾经适应的那个环境，而只有在这样的环境下我们的身体状态才是正常的。换句话说，锻炼是原始人身体活动的替代。这也是我们需要做的、从而能使身在21世纪的我们身心正常的一件事。何况只是身心正常，对有追求的人来说还是远远不够的。

一个运动员决定开始某种力量训练计划的动机，可能只是想要参加一项需要这种力量的团队运动，或者出于某些更私人的原因。但很多人只是感觉自己的力量不够，或者觉得自己还可以继续提高，而不是想要加入某个运动队。这本书就是为这类人创作的。

\section{为什么选择杠铃？}

力量训练和人类文明一样，具有悠久的历史。古希腊关于大力士米罗（Milo）的传说证明了自古以来人们就热衷于通过力量训练让身体变强，并且让人们了解了身体变强所需要的过程。据说米罗每天都背着一只牛犊行走，随着牛犊越长越大，他也变得日益强壮。在几千年前，人们就已经认识到发展力量是一个渐进性的过程，但直到现在，人们才通过技术解决了如何才能有效地进行渐进式阻力训练的问题。

杠铃是人们最早发明的进行阻力训练的器材之一——一根长的金属杆两头挂上某种形式的重物。最早的杠铃使用球体作为配重，通过它们可以让杠铃保持平衡，并且能够通过在球体中加入沙子或者铅弹来增加杠铃的重量。戴维·威洛比（David Willoughby）1970年创作的一本畅销书——《超级运动员》（\textit{The Super Athlete}）详细说明了举重和相应器材的历史。

但是威洛比先生没有预见的是，举重器材的发展在20世纪70年代中期发生了剧烈的改变。一位叫作阿瑟·琼斯（Arthur Jones）的先生发明了一种对阻力训练产生革命性改变的器械。不幸的是，并非所有的革命性改变都能够有真正的贡献。琼斯运用了“变化阻力法则”——这个法则利用了四肢在整个动作幅度内的某一位置比其他位置更强的原理——设计出了锻炼不同肢体或者身体部位的机器：一个滑轮与连接重物的铁链组合在一起，在运动过程中给关节提供不同的阻力。这种机器被设计出来的目的是为了训练者能够按照一定的次序来锻炼，一组接着一组，不需要组间休息，因为进行连续练习的是不同的身体部位。这个机器的核心思想（从商业角度来讲）是，如果有足够的机器——每个机器分别锻炼我们的一个身体部位——然后将其全部集合起来组成一个机器组的话，那么我们身体的每个部位都能被锻炼到。这些机器做工十分精细，外观漂亮，很快就成了大多数健身房的标配，这就是价格昂贵的12站式诺德士（Nautilus）牌健身器。

训练机器其实并没有多么新奇。大多数高中都会有一个“环球角斗士”（Universal Gladiator）多站式装置，而腿部屈伸机和锻炼背阔肌的机器对于每个做过重量训练的人来说都不陌生，区别在于新机器背后的市场运作。诺德士公司着重宣传了完整的循环训练会起到怎样明显的效果，而这在以前从未被重视过。他们推出了一系列展示训练者训练前后变化的广告，广告中的凯西·维亚托（Casey Viator）在只使用诺德士牌机器训练后肌肉增长明显。但有一点在广告中被忽略了——作为一名有经验的健美运动员，使用诺德士牌健身器只是让维亚托先生的肌肉恢复到了之前使用传统训练方法就已经达到的效果而已。

琼斯甚至申明，使用诺德士牌健身器所增长的身体力量能转移到诸如奥林匹克举重这样的复合型动作中去，而不需要借助大重量来专门练习这些动作。这种说法与传统的训练理论和训练者的实际训练经验完全相悖。但是其说法还是被越来越多的人认可，因此诺德士公司取得了巨大的商业成功：他们的机器作为现代商业健身器械的标准风靡全球。

诺德士牌健身器之所以如此成功，是因为它的出现让健身房行业给大众提供了一种从未有过的体验。在诺德士牌健身器问世之前，如果健身房的会员需要的训练强度远远大于“环球”机器所能允许的上限的话，那么他不得不学习使用杠铃去训练。必须有人教他使用杠铃，而且还必须有人教会健身房的员工如何教他使用杠铃。这种专业的教育在以前和现在都是很费时间，并且不太普及的。但是如果有了诺德士牌健身器，即使拿着最低薪水的员工也能很快掌握其使用流程。很显然，健身房经营者无须在员工培训上投入太多，就能够为会员们提供一整套全身训练方案。此外，完成整套流程大约只需30分钟，这样就大大减少了会员在俱乐部滞留的时间，从而增加了客流量并且能够使营业额最大化。诺德士牌健身器真正使现代健身房的出现成为可能。

但是问题在于，以机器为基础的训练并不像广告中描述的那样有效。完成一个流程的训练几乎不可能增加肌肉的含量。尝试着这样做的人会发现，努力训练了几个月而肌肉量并没有明显的增加。但当他们改用杠铃训练时，奇迹就会发生：在一个星期内，他们增长的肌肉量会比他们在机器上奋斗了几个月增长的还要多。

把不同的身体部位孤立起来，各自独立地进行训练是使用机器训练没有效果的根源所在。而杠铃训练之所以有效果的原因在于，训练者可以同时锻炼全身肌肉。杠铃优于其他任何增长力量的健身器材。人体是作为一个完整的系统来运作的——它以这样的方式运作，它也希望能以同样的方式进行训练。人体不喜欢被分成各个部分，然后各自孤立开来训练，因为训练获得的力量不是以这种方式被使用的。获取力量的模式必须与使用力量的模式相同。神经系统连接、控制着肌肉，当你获取力量的方式没有和实际使用力量的方式对应起来的话，你就没有考虑到这种连接。不幸的事实是，神经肌肉的连接是专一的——锻炼计划必须遵守这个原则，就像遵守重力法则一样。

杠铃和基于杠铃的主要练习远远优于其他的训练器材。\textbf{正确实施的、全动作幅度的杠铃练习，实际上是人类骨骼、肌肉系统在负重下的功能性表达。}杠铃练习由每个人特定的运动模式来控制，并由训练者肢体的长度、肌肉的连接位置、力量水平、柔韧性和神经肌肉连接效能来精细调整。在一个杠铃动作中，所有参与的肌肉是自然达到平衡状态的，因为所有参与的肌肉都贡献了它们在人体结构上特定的那一部分工作量。肌肉移动骨骼之间的关节，骨骼则把力传递给杠铃：这样的运作方式是符合人体运动系统的规律的——当人体系统根据自己的设计运转时，它会达到最优状态，所以我们的训练应该遵循其设计。杠铃移动的方式恰好符合人体设计的动作模式，因为动作的每一方面都由身体自己决定。

而在另一方面，机器强迫身体根据它的设计来移动相应的重量，这在很大程度上限制了练习者所获得的能力，使其不能很好地满足某一特定运动的需求。比如说，一个人不可能在任何运动中孤立地使用股四头肌，而把腘绳肌撇开，这种运动方式只存在于专为此目的而设计的机器中——没有任何自然的运动是这样的。股四头肌和腘绳肌总是共同发挥作用，同时从两边来平衡膝关节所受的力。既然它们总是一起运作，为什么要把它们分开来训练呢？仅仅是因为有人发明了一种能让我们这样做的机器吗？

就算机器允许多关节同时运动，也达不到理想的训练状态，因为人体在空间中的运动模式是由机器决定的，而不是基于人体固有的生物力学特点。但是杠铃允许训练者在运动过程中进行微调，从而适应训练者个人的身体结构特点。

不仅如此，在训练时杠铃还要求训练者进行所有必要的调整，以在重量移动的过程中保持控制。这个方面的重要性再怎么强调都不为过——对杠铃的控制、对训练者平衡和协调能力的需求，这些是杠铃训练所独有的，而在以机器为基础的训练中是完全见不到的。因为负重运动的每一个方面都要由训练者自己控制，这样每一个方面都会得到锻炼。

杠铃训练还有其他好处。本书中描述的所有练习都涉及不同程度的骨骼负载。毕竟，杠铃的重量最终是由骨骼来支撑的。骨骼是有生命的、能够对外界的压力做出反应的组织，就像肌肉、韧带、肌腱、皮肤、神经和大脑一样。它也会像其他组织一样适应外界的刺激，并在经过更大重量的训练后变得密度更大、更结实。从这方面来说，杠铃训练对老年人和女性来说是很重要的，因为骨密度是影响他们健康的一个主要因素。

使用杠铃是非常经济实惠的。购买一套任一品牌的现代训练机器的钱，都可以用来组建五六个非常实用的举重室——你可以在里面做几百种不同的练习。就算对你来说钱不是问题，你也应当考虑实用性。在一个公共机构中，每花一元钱能够在一定时间内让多少人得到训练，也许是你决定购买何种器械的重要考量。正确的选择也许会直接影响你的训练体验。

杠铃训练的唯一问题是相当多的人不知道如何正确地使用杠铃。学习渠道的缺失“名正言顺”地阻止了很多人使用杠铃进行训练。我将通过这本书来尝试解决这个问题。本书中教授的杠铃训练方法已经在健身行业中发展了三十多年，这种训练方法的一小部分内容仍然掌握在这些人手中——他们追求成果、能够坦诚对待真正有用的训练方法，并且尊重久经考验的生物科学法则。我希望它也能对你们有用——就像对曾经的我那样。

\begin{quote}
    这套奥林匹克杠铃片组是在威奇托福尔斯（Wichita Falls）市区的基督教青年会找到的。它有将近50年的历史，被数以千计的男男女女使用过。在这其中有比尔·斯塔尔——著名的力量训练教练、奥林匹克举重运动员，也是最初的力量举比赛的竞赛者之一。比尔是霍夫曼（Hoffman）《力量和健康》（\textit{Strength and Health}）杂志和乔·韦德（Joe Weider）《肌肉》（\textit{Muscle}）杂志的编辑。他是第一批专业级别的全职力量教练的一员，还曾在很多的国家级、世界级的奥林匹克举重队中担任过教练。他是介绍杠铃训练方面最多产的作家之一，在超过50年的时间中出版了很多书籍，发表了很多文章。他的训练伙伴和他训练过的运动员成就非凡，这使他的影响力延续至今。他最初的举重训练就是使用这套杠铃片组完成的。

    \begin{flushright}
        ——摘自德克萨斯州威奇托福尔斯市运动俱乐部的比尔·斯塔尔纪念碑
    \end{flushright}
\end{quote}

\begin{figure}[htbp] % htbp 控制图片排版位置优先顺序
    \centering
    \includegraphics[width=0.8\textwidth]{Image00002.jpg} % 插入图片，宽度设为版心的 80%
\end{figure}
